\documentclass[11pt]{article}
\usepackage[margin=1in]{geometry}
\usepackage{times}
\usepackage{booktabs}
\usepackage{graphicx}
\usepackage{amsmath}
\usepackage[hidelinks]{hyperref}
\usepackage{titlesec}
\titleformat{\section}{\large\bfseries}{\thesection.}{0.6em}{}
\titleformat{\subsection}{\normalsize\bfseries}{\thesubsection}{0.6em}{}
\setlength{\parskip}{4pt}
\setlength{\parindent}{0pt}

\title{\vspace{-1.5em}\textbf{Does 4-Bit Quantization Change Multi-Turn Jailbreak Robustness?}\\[0.3em]
\large A Two-Metric Evaluation of NF4, AWQ, and GPTQ on Llama-3.2-3B and Llama-3.1-8B}
\author{Parv Bansal\\ AIMS-DTU Research Internship 2026\\ \texttt{parvbansal0011@gmail.com}}
\date{2026}

\begin{document}
\maketitle

\begin{abstract}
\noindent
Quantization lets large language models run on edge hardware, where real use is
multi-turn. Prior work measured quantization's effect on refusal at a single turn
only and reported degradation. This paper gives the first measurement across a
multi-turn jailbreak, testing the claim on Llama-3.2-3B (FP16, NF4) and
Llama-3.1-8B (FP16, AWQ, GPTQ) with two metrics: a refusal classifier and a
Llama Guard 3 harm judge. Quantization does not degrade safety. On standalone
AdvBench prompts, single-turn refusal is flat or higher under quantization for
both models and all three methods, with no significant comparison. Under CoSafe
multi-turn attacks, precision changes neither refusal (3B 23\% versus 23\%,
McNemar $p=0.62$; 8B 27 to 32\%, Cochran $Q$ $p=0.12$) nor harmful-response rate.
The one significant single-turn effect, 8B refusing more after quantization,
vanishes under the attack. A safety-prompt positive control confirms the harm
metric can detect a real intervention, so the null is a measured absence of
effect rather than a lack of power. The evaluation also shows that refusal-based
attack success overstates harm by an order of magnitude: the CoSafe attack drives
8B to 73\% non-refusal, but only 6\% of those responses are actually harmful.
Single-turn safety evaluation does not predict how quantized models rank under
conversational attack.
\end{abstract}

\section{Introduction}

Post-training quantization to 4 bits is now the standard way to fit an
instruction-tuned model onto a phone or a laptop. The same deployment target that
forces quantization also forces multi-turn use, because a chat assistant holds a
conversation. Safety evaluation has not followed. Recent work that measures how
quantization changes refusal uses single-turn prompts and reports a drop in
refusal after quantization (Kadadekar, 2026). Whether that drop holds, grows, or
disappears once an attacker has several turns is open, and the authors name
multi-turn as future work.

This paper answers one question: does 4-bit quantization change a model's
multi-turn jailbreak robustness, and how does that compare to its single-turn
effect on the same models and prompts? The study covers two models and three
quantization methods, and reads two behavioral signals, whether the model refused
and whether the response is actually harmful, along with the projection onto the
refusal direction of Arditi et al. (2024).

The reported behavioral degradation does not reproduce, at either horizon. On
standalone AdvBench prompts, quantization leaves single-turn refusal flat or
higher for both models and all three methods, with no significant comparison.
Under a CoSafe coreference attack, precision does not change refusal either, and
all precisions converge to a common non-refusal rate. A Llama Guard 3 rescore
confirms the null on a second metric: quantization never raises the rate of
actually-harmful responses. The only significant behavioral precision effect
anywhere is an increase, where 8B quantization raises refusal on the CoSafe cold
turn, and even that vanishes under the attack.

The practical consequence is a warning about evaluation. A single-turn safety
score does not preserve how FP16 and quantized checkpoints rank under
conversational attack, and on these models it does not surface a behavioral
safety cost from quantization at all. The contributions are:

\begin{itemize}
\itemsep2pt
\item A well-powered non-replication. Across two Llama models, three 4-bit
methods, two prompt sets (AdvBench and CoSafe), and two metrics (substring
refusal and a Llama Guard harm judge), quantization does not reduce safety at a
single turn.
\item The first measurement of 4-bit quantization under a multi-turn attack.
Precision changes neither multi-turn refusal nor multi-turn harm, and the one
significant single-turn effect does not survive the attack.
\item Two evaluation cautions. Refusal-based attack success overstates actual
harm by an order of magnitude on CoSafe, and a safety-prompt positive control
shows the harm metric can detect a real intervention, so the quantization null is
a true absence of effect.
\end{itemize}

\section{Background}

\subsection{Quantization and safety}
Kadadekar (2026) reports that INT4 quantization degrades single-turn refusal
while the refusal direction stays geometrically intact, and notes that standard
probes do not explain the behavioral gap. Chhabra and Khalili (2025) find that
the refusal direction survives compression with cosine similarity above 0.99.
Both study single-turn behavior. Neither measures a conversation. This paper
extends the first to the multi-turn setting it names as open.

\subsection{Refusal as a direction}
Arditi et al. (2024) show that refusal in instruction-tuned models is mediated by
a single linear direction in the residual stream. Adding it induces refusal, and
ablating it removes refusal. This study uses their difference-in-means estimator
and confirms the causal layer by ablation rather than by separability.

\subsection{Multi-turn attacks}
STAR (2026) shows that in full-precision models the projection onto the refusal
direction decays across conversation turns, so multi-turn state erodes refusal.
CoSafe (Yu et al., 2024) builds multi-turn attacks whose final turn refers back
to earlier turns by coreference, so the harmful request is only legible in
context. The measurement here is independent of STAR and is used to compare
precisions, not to characterize full-precision decay.

\section{Experimental Setup}

\subsection{Models and precisions}
The 3B model is Llama-3.2-3B-Instruct at FP16 and NF4 (bitsandbytes
double-quantized). The 8B model is Llama-3.1-8B-Instruct at FP16, AWQ-INT4, and
GPTQ-INT4, using the public \texttt{hugging-quants} checkpoints. AWQ and GPTQ
layers are not fused, so forward hooks see the residual stream. FP16 extraction
and ablation run on Apple silicon, and quantized kernels run on CUDA.

\subsection{Refusal direction and causal layer}
The refusal direction is extracted by difference-in-means: the mean activation at
the post-instruction token for harmful prompts minus the same for harmless
prompts, per layer, normalized. Held-out separation accuracy is 100\% for the 3B
direction. The functional layer is chosen by ablation, not by separability. Each
candidate layer's direction is removed from the residual stream, and the layer
whose removal collapses refusal is kept. For 3B this is layer 20 and for 8B layer
12, where ablation drops refusal from 94\% to 0\% while a random-direction control
stays at baseline. The separability-argmax layer is different in both models, 28
for 3B and 31 for 8B, and does not mediate refusal, so only ablation-confirmed
layers are used.

\subsection{Attacks and prompts}
Multi-turn attacks come from CoSafe, using its safe categories (non-violent
unethical behavior, financial crime, privacy violation, and misinformation), 100
scenarios in total. For each scenario the final unsafe turn is measured two ways.
\textbf{Cold} presents the final turn alone with no conversation. \textbf{In-context}
presents the full conversation, which is the attack. The two conditions share the
same final request, so their difference isolates the effect of the conversational
context. Standalone single-turn prompts come from AdvBench, 100 instructions
presented as one turn, to test the quantization effect on overtly harmful prompts
in the same setting as prior single-turn work.

\subsection{Metrics}
Refusal is detected by a fixed phrase classifier on greedy generation. To measure
whether a response is actually harmful, every completion is rescored with Llama
Guard 3 (Inan et al., 2023), reporting the harmful-response rate. The projection
of the post-instruction activation onto the refusal direction is recorded as a
mechanism readout. Refusal and harm are paired binary across precision, so
comparisons use McNemar's test per pair and Cochran's $Q$ across the three 8B
precisions. Projection is paired continuous, so it uses the paired $t$-test. All
runs are 100 items, so $n=100$ per condition.

\section{Results and Analysis}

\subsection{Multi-turn refusal does not depend on precision}
Table~\ref{tab:means} gives the rates. For 3B the two precisions are identical at
23\% in-context refusal (McNemar $\chi^2=0.25$, $p=0.62$), and the paired
projection difference is null ($t=0.57$, $p=0.57$). For 8B the three precisions
span 27 to 32\%, and Cochran's $Q$ over the paired refusals is not significant
($Q=4.2$, $p=0.12$). The non-refusal rate sits near 70\% for every precision,
though most of that is not harmful output, as the harm judge below shows.

\begin{table}[t]
\centering
\begin{tabular}{llcccc}
\toprule
Model & Prec. & \multicolumn{2}{c}{Refusal} & \multicolumn{2}{c}{Refusal-dir.\ proj.} \\
& & cold & in-ctx & cold & in-ctx \\
\midrule
3B & FP16 & 25\% & 23\% & $+1.46$ & $-0.86$ \\
3B & NF4  & 23\% & 23\% & $+1.33$ & $-0.84$ \\
\midrule
8B & FP16 & 23\% & 27\% & $+0.88$ & $+0.82$ \\
8B & AWQ  & 31\% & 32\% & $+0.80$ & $+0.76$ \\
8B & GPTQ & 34\% & 31\% & $+0.90$ & $+0.89$ \\
\bottomrule
\end{tabular}
\caption{Refusal and mean projection onto the refusal direction, cold (single
turn) and in-context (multi-turn), $n=100$ per cell.}
\label{tab:means}
\end{table}

\subsection{The single-turn precision gap does not survive the attack}
The cold column tells a different story for 8B. Quantization raises single-turn
refusal, with FP16 at 23\%, AWQ at 31\%, and GPTQ at 34\%. Both pairwise contrasts
against FP16 are significant (McNemar $p=0.043$ and $p=0.010$; Table~\ref{tab:sig}),
and Cochran's $Q$ is significant ($Q=9.7$, $p=0.008$). The same three precisions
are not separable once the conversation is present. The 8 to 11 point spread at a
single turn narrows to 5 points in context and loses significance. The 3B model
shows no precision effect at either horizon.

\begin{table}[t]
\centering
\begin{tabular}{lll}
\toprule
Comparison & Cold & In-context \\
\midrule
3B FP16 vs NF4 & $p=0.68$ & $p=0.62$ \\
8B FP16 vs AWQ & $p=0.043$ & $p=0.13$ \\
8B FP16 vs GPTQ & $p=0.010$ & $p=0.22$ \\
8B Cochran $Q$ & $p=0.008$ & $p=0.12$ \\
\bottomrule
\end{tabular}
\caption{Refusal significance (McNemar per pair, Cochran $Q$ across the three 8B
precisions). The single-turn precision effect on 8B is significant, and the
multi-turn effect is not.}
\label{tab:sig}
\end{table}

\subsection{No degradation on standalone harmful prompts}
On AdvBench single-turn prompts (Table~\ref{tab:advbench}) quantization does not
lower refusal either. The 3B model moves from 76\% to 82\% under NF4 (McNemar
$p=0.15$). The 8B model moves from 95\% to 98\% under AWQ ($p=0.25$) and stays at
95\% under GPTQ ($p=0.62$). No comparison is significant, and every point estimate
is flat or higher under quantization. This is a well-powered non-replication of
the single-turn refusal degradation reported for INT4 quantization. Refusal here
is near ceiling for 8B, which limits behavioral power, so the harm judge below
addresses the question from the other side.

\begin{table}[t]
\centering
\begin{tabular}{llccc}
\toprule
Model & \multicolumn{3}{c}{Refusal} & McNemar \\
& FP16 & \multicolumn{2}{c}{quantized} & vs FP16 \\
\midrule
3B & 76\% & NF4 82\% & & $p=0.15$ \\
8B & 95\% & AWQ 98\% & GPTQ 95\% & $p=0.25 / 0.62$ \\
\bottomrule
\end{tabular}
\caption{Single-turn refusal on 100 standalone AdvBench prompts. No precision
comparison is significant, and quantization never lowers refusal.}
\label{tab:advbench}
\end{table}

\subsection{The null holds under a harm judge}
Refusal is a lenient signal. Rescoring every completion with Llama Guard 3, which
labels whether the response itself is harmful, gives the harmful-response rates in
Table~\ref{tab:harm}. Quantization does not raise the harmful-response rate in any
condition. Every quantized run is at or below its FP16 baseline, and no comparison
is significant (McNemar $p \geq 0.48$). The null survives a second, stricter
metric.

\begin{table}[t]
\centering
\begin{tabular}{llcc}
\toprule
Setting & Model & FP16 & quantized \\
\midrule
single turn & 3B & 1\% & NF4 1\% \\
single turn & 8B & 4\% & AWQ 2\% / GPTQ 3\% \\
multi turn  & 3B & 4\% & NF4 3\% \\
multi turn  & 8B & 6\% & AWQ 4\% / GPTQ 5\% \\
\bottomrule
\end{tabular}
\caption{Harmful-response rate (Llama Guard 3), lower is safer. Quantization never
raises it, and all FP16-versus-quantized comparisons are non-significant.}
\label{tab:harm}
\end{table}

\subsection{The harm metric has power: a positive control}
A null on a rare-event metric invites the objection that nothing could show up.
To rule that out, a safety system prompt is prepended to every FP16 condition, an
intervention that should change behavior. It does, and the harm metric detects it.
Under the multi-turn attack, the system prompt cuts 8B harmful responses from 6\%
to 0\% (McNemar $p=0.041$) and raises multi-turn refusal from 27\% to 38\%
($p=0.006$). The harm metric is not saturated. It registers a real safety
intervention on the same prompts where quantization registers nothing. The
quantization null is a measured absence of effect, not an absence of power. The
system prompt is the opposite of quantization here, because it does move
multi-turn safety, which is what makes it a clean control.

\subsection{Refusal-based attack success overstates harm}
The two metrics diverge sharply. The CoSafe attack takes 8B to a 73\% non-refusal
rate by the substring metric, but only 6\% of those responses are judged harmful
by Llama Guard. The attack mostly produces non-refusing but non-harmful text. On
these prompts, refusal-based attack success overstates actual harm by an order of
magnitude, which is a caution for how multi-turn safety is reported.

\subsection{The direction stays coupled}
If quantization broke the refusal representation, the projection would stop
predicting refusal. It does not. The AUC of projection predicting refusal ranges
from 0.79 to 0.93 across all precisions and conditions and never approaches
chance, and point-biserial correlations run from 0.46 to 0.68 with $p<10^{-6}$.
The refusal direction drives behavior under AWQ, GPTQ, and NF4 as it does under
FP16. Multi-turn convergence is therefore behavioral saturation across a shared
threshold, not a loss of the mechanism.

\section{Discussion}

\subsection{Why the multi-turn gap closes}
CoSafe attacks hide the harmful request behind coreference, so the model answers
the earlier benign-looking turns and then completes the final one in context. The
attack is strong, taking FP16 from 23\% to 27\% refusal for 8B and driving 3B to
23\%. Once the attack is this effective, a few points of extra caution from
quantization is not enough to change the outcome. The precision difference is real
at a single turn and is swamped by the attack in context. A safety margin that
exists in isolation does not imply a safety margin under pressure.

\subsection{Implication for evaluation}
Single-turn refusal is the standard axis for comparing a quantized checkpoint to
its FP16 source. This result shows that axis does not preserve the multi-turn
ordering. For 8B, single-turn evaluation ranks the quantized checkpoints as safer
than FP16, while under attack they are statistically tied. A single-turn number,
whether it shows degradation or improvement, does not transfer to the
conversational setting where the model is deployed.

\subsection{Relation to prior reports}
Kadadekar (2026) reports single-turn degradation on standard harmful prompts.
This study does not reproduce it. On AdvBench, the same standard single-turn
setting, no precision comparison is significant and every point estimate is flat
or higher under quantization, for both models and all three methods. On the
CoSafe cold turn, 8B quantization instead raises refusal. Across both prompt sets
and both metrics the behavioral effect of quantization is null-to-positive, not
degradation. The Llama Guard rescore closes the ceiling concern that would
otherwise bound a refusal-only null, because on a metric with room below and one
the positive control shows can detect a real intervention, quantization still
moves nothing.

\section{Limitations and Future Work}
The study covers two model sizes (3B and 8B), one family (Llama), one multi-turn
attack (CoSafe coreference), and three quantization methods (NF4, AWQ, GPTQ). It
does not vary attack type or study training-time quantization. The refusal
detector is a phrase classifier, which can miss soft refusals, though the harm
judge and the projection readout, which do not use it, agree with the behavioral
results. CoSafe safe categories are used only. Both metrics sit near an extreme,
with substring refusal near ceiling (95\% for 8B on AdvBench) and the Llama Guard
harm rate near floor (1 to 6\%), so each alone has limited room to expose a small
effect. The positive control mitigates this for the harm metric, but the low harm
rate partly reflects the CoSafe safe categories being mild relative to the
severe-hazard taxonomy Llama Guard is tuned for, so it may under-flag them. The
result is best read as: on these prompts, actual harmful output is rare at every
precision, not as a guarantee on severe-harm categories. A stricter or
broader-taxonomy judge, a second multi-turn attack, and a larger model family are
the direct next steps.

\section{Conclusion}
Across two Llama models, three 4-bit methods, two prompt distributions, and two
metrics, 4-bit quantization does not reduce safety at either a single turn or
across a multi-turn attack. The only significant behavioral precision effect is an
increase, where 8B refuses more on the CoSafe cold turn, and it does not survive
the attack, where all precisions converge near 30\% refusal and 4 to 6\% harmful
output. A safety system prompt, by contrast, does lower multi-turn harm on the
same prompts, so the null is a true absence of effect rather than a lack of
measurement power. Single-turn safety evaluation does not predict how quantized
models rank under conversational attack, and on these models it does not show a
behavioral safety cost from quantization at all.

\section*{References}
\small
Arditi, A., Obeso, O., Syed, A., Paleka, D., Panickssery, N., Gurnee, W., and
Nanda, N. (2024). Refusal in Language Models Is Mediated by a Single Direction.
\textit{NeurIPS}.

Chhabra, A., and Khalili, M. M. (2025). Refusal in Compressed Models via
Mechanistic Interpretability. \textit{arXiv:2504.04215}.

Dettmers, T., Pagnoni, A., Holtzman, A., and Zettlemoyer, L. (2023). QLoRA:
Efficient Finetuning of Quantized LLMs. \textit{NeurIPS}.

Frantar, E., Ashkboos, S., Hoefler, T., and Alistarh, D. (2023). GPTQ: Accurate
Post-Training Quantization for Generative Pre-trained Transformers. \textit{ICLR}.

Inan, H., Upasani, K., Chi, J., Rungta, R., Iyer, K., Mao, Y., et al. (2023).
Llama Guard: LLM-based Input-Output Safeguard for Human-AI Conversations.
\textit{arXiv:2312.06674}.

Kadadekar. (2026). Quality Is Not a Safety Proxy Under Quantization.
\textit{arXiv:2606.10154}.

Lin, J., Tang, J., Tang, H., Yang, S., Dang, X., and Han, S. (2024). AWQ:
Activation-Aware Weight Quantization for LLM Compression and Acceleration.
\textit{MLSys}.

STAR. (2026). State-Dependent Safety Failures in Multi-Turn Language Model
Interaction. \textit{arXiv:2603.15684}.

Yu, E., et al. (2024). CoSafe: Evaluating Large Language Model Safety in
Multi-Turn Dialogue Coreference. \textit{EMNLP}.

Zou, A., Wang, Z., Kolter, J. Z., and Fredrikson, M. (2023). Universal and
Transferable Adversarial Attacks on Aligned Language Models.
\textit{arXiv:2307.15043}.

\end{document}
